% 06-networking.tex
\chapter{Networking}
\label{ch:networking}

\section{Network Architecture}
\label{sec:network-architecture}

Doki's networking stack provides bridge networks, CNI plugin support,
port mapping, an internal DNS server, and DokiLink-Lite peer-to-peer
mesh networking. The stack auto-detects available capabilities and
adjusts behavior accordingly.

\begin{figure}[htbp]
  \centering
  \begin{tikzpicture}[
    node distance=0.5cm,
    every node/.style={font=\sffamily\footnotesize},
  ]
    % Bridge network (left side)
    \node[comp, minimum width=3.5cm] (bridge) {doki0 Bridge\\{\scriptsize\color{textgray} 10.0.0.0/24}};
    \node[compSm, below left=0.5cm and 0.2cm of bridge] (c1) {Container 1\\{\scriptsize\color{textgray} 10.0.0.2}};
    \node[compSm, below=0.5cm of bridge] (c2) {Container 2\\{\scriptsize\color{textgray} 10.0.0.3}};
    \node[compSm, below right=0.5cm and 0.2cm of bridge] (c3) {Container 3\\{\scriptsize\color{textgray} 10.0.0.4}};

    \draw[arrDash] (bridge) -- (c1);
    \draw[arrDash] (bridge) -- (c2);
    \draw[arrDash] (bridge) -- (c3);

    % DNS (middle)
    \node[compAlt, right=1cm of bridge] (dns) {DNS Server\\{\scriptsize\color{textgray} 127.0.0.11:8053}};

    % DokiLink (right)
    \node[compTeal, right=1cm of dns] (mesh) {DokiLink\\{\scriptsize\color{textgray} TLS 1.3}};

    \draw[arr] (bridge) -- (dns);
    \draw[arr] (dns) -- (mesh);

    % Region
    \node[region, fit=(bridge)(c1)(c3), label={[annotBold]above:Local Bridge Network}] {};
  \end{tikzpicture}
  \caption{Doki networking components: bridge network, DNS server,
    and DokiLink mesh.}
  \label{fig:networking-overview}
\end{figure}

\section{Bridge Networking}
\label{sec:bridge}

The default network type is a Linux bridge (\texttt{doki0}) with NAT,
DNS, and port mapping. When running on Android without root, Doki
falls back to socat-based port forwarding.

\begin{figure}[htbp]
  \centering
  \begin{tikzpicture}[
    node distance=0.6cm,
    every node/.style={font=\sffamily\footnotesize},
  ]
    % Host network
    \node[rectangle, rounded corners=3pt, draw=nodeblue!50, fill=fillblue!50, minimum width=3cm, minimum height=0.7cm, text=textdark, align=center, inner sep=4pt] (host) {Host Network};

    % Bridge
    \node[rectangle, rounded corners=3pt, draw=nodeorange!60, fill=fillorange!30, minimum width=3.5cm, minimum height=0.7cm, text=textdark, align=center, inner sep=4pt, below=0.8cm of host] (bridge) {doki0 Bridge\\{\scriptsize\color{textgray} 10.0.0.0/24}};

    % Containers
    \node[rectangle, rounded corners=2pt, draw=nodelteal!50, fill=fillteal!30, minimum width=1.8cm, minimum height=0.6cm, text=textdark, align=center, inner sep=3pt, below left=0.8cm and 0.3cm of bridge] (c1) {Container 1\\{\scriptsize\color{textgray} .0.2}};
    \node[rectangle, rounded corners=2pt, draw=nodelteal!50, fill=fillteal!30, minimum width=1.8cm, minimum height=0.6cm, text=textdark, align=center, inner sep=3pt, below=0.8cm of bridge] (c2) {Container 2\\{\scriptsize\color{textgray} .0.3}};
    \node[rectangle, rounded corners=2pt, draw=nodelteal!50, fill=fillteal!30, minimum width=1.8cm, minimum height=0.6cm, text=textdark, align=center, inner sep=3pt, below right=0.8cm and 0.3cm of bridge] (c3) {Container 3\\{\scriptsize\color{textgray} .0.4}};

    % Port mapping
    \node[rectangle, rounded corners=2pt, draw=nodepurple!50, fill=fillpurple!30, minimum width=2cm, minimum height=0.6cm, text=textdark, align=center, inner sep=3pt, left=1.5cm of c1] (port) {Port Mapping\\{\scriptsize\color{textgray} 8080$\to$80}};

    % NAT
    \node[rectangle, rounded corners=2pt, draw=nodecoral!50, fill=fillcoral!30, minimum width=1.5cm, minimum height=0.6cm, text=textdark, align=center, inner sep=3pt, right=1.5cm of c3] (nat) {NAT\\{\scriptsize\color{textgray} masquerade}};

    % Arrows
    \draw[arr] (host) -- (bridge);
    \draw[arrDash] (bridge) -- (c1);
    \draw[arrDash] (bridge) -- (c2);
    \draw[arrDash] (bridge) -- (c3);
    \draw[arr] (port) -- (c1);
    \draw[arr] (c3) -- (nat);

    % Labels
    \node[annot, above right=0.1cm and 0.2cm of host.north east] {\textbf{Host}};
    \node[annot, below=0.2cm of nat] {to internet};
  \end{tikzpicture}
  \caption{Network topology showing the bridge, containers, port mapping, and NAT.}
  \label{fig:network-topology}
\end{figure}

\subsection{Port Mapping}
\label{sec:port-mapping}

Port mapping rules are applied to both the PREROUTING chain (for
external traffic) and the OUTPUT chain (for locally-originated
traffic). This ensures that port mappings work correctly regardless
of where the connection originates.

The nftables implementation constructs rules in the following format:

\begin{codebox}[title={Port mapping rule structure}]{text}
iptables -t nat -A DOKI -p tcp --dport [host\_port] \\
  -j DNAT --to-destination [container\_ip]:[container\_port]
\end{codebox}

\section{Internal DNS}
\label{sec:internal-dns}

Doki runs an internal DNS server that resolves container names to IP
addresses. The DNS server supports A, AAAA, and PTR records with an
LRU cache of 1024 entries and a 5-minute TTL.

\subsection{Port Selection}
\label{sec:dns-port}

On Android, the DNS server listens on \texttt{127.0.0.11:8053}
instead of the standard port 53, which is blocked by SELinux policies.
The \texttt{resolv.conf} generation uses
\texttt{net.SplitHostPort()} to strip the port number, since
\texttt{resolv.conf} does not support port specifications.

\subsection{DNS Resolution Flow}
\label{sec:dns-resolution}

When a container performs a DNS lookup, the following steps occur:

\begin{figure}[htbp]
  \centering
  \begin{tikzpicture}[
    node distance=0.5cm,
    every node/.style={font=\sffamily\footnotesize},
    dnsStep/.style={
      rectangle, rounded corners=2pt,
      draw=nodeblue!60, fill=fillblue,
      minimum width=2.5cm, minimum height=0.7cm,
      text=textdark, align=center, inner sep=4pt
    }
  ]
    \node[dnsStep] (query) {1. Query\\{\scriptsize\color{textgray} from container}};
    \node[dnsStep, right=0.8cm of query] (cache) {2. Cache\\{\scriptsize\color{textgray} lookup}};
    \node[dnsStep, right=0.8cm of cache] (miss) {3. Cache\\{\scriptsize\color{textgray} miss}};
    \node[dnsStep, right=0.8cm of miss] (upstream) {4. Upstream\\{\scriptsize\color{textgray} DNS query}};
    \node[dnsStep, right=0.8cm of upstream] (tcp) {5. TCP\\{\scriptsize\color{textgray} retry (RFC 5966)}};
    \node[dnsStep, right=0.8cm of tcp] (response) {6. Cache +\\{\scriptsize\color{textgray} return}};

    \draw[arr] (query) -- (cache);
    \draw[arr] (cache) -- (miss);
    \draw[arr] (miss) -- (upstream);
    \draw[arr] (upstream) -- (tcp);
    \draw[arr] (tcp) -- (response);

    % Cache hit arrow
    \draw[arrDash, nodeorange] (cache.south) -- ++(0,-0.5)
      node[below, font=\sffamily\tiny, text=nodeorange] {cache hit}
      -- ++(2.5,0) |- (response.south);
  \end{tikzpicture}
  \caption{DNS resolution flow with cache and TCP fallback.}
  \label{fig:dns-flow}
\end{figure}

When a container performs a DNS lookup, the following steps occur:

\begin{enumerate}[leftmargin=*, itemsep=3pt]
  \item The query arrives at \texttt{127.0.0.11:8053}
  \item The server checks the LRU cache for a recent response
  \item On cache miss, the server queries upstream DNS servers
  \item If the response has the TC (truncated) bit set, the server
        retries using TCP (RFC 5966)
  \item The response is cached and returned to the container
\end{enumerate}

\section{DokiLink-Lite Mesh Networking}
\label{sec:dokilink}

DokiLink-Lite provides peer-to-peer encrypted communication between
Doki instances. It uses a gossip protocol for peer discovery and
supports three encryption layers.

\subsection{Encryption Layers}
\label{sec:encryption-layers}

\begin{description}[leftmargin=*, style=nextline, font=\sffamily\bfseries]
  \item[L1: TLS 1.3] (Default)
    Each link certificate is signed by a per-install ECDSA P-256
    Certificate Authority. Certificates include SAN DNS names for
    identity verification. The CA and keypair are generated on first
    start and stored at \texttt{\$DOKI\_DATA\_DIR/keys/}.

  \item[L2: NaCl Secretbox] (Opt-in)
    When \texttt{DOKI\_LINK\_PAYLOAD\_ENC=1} is set, a 32-byte
    encryption key is derived from both peers' Ed25519 public keys
    using SHA-256. The key is order-sensitive, ensuring both peers
    derive the same key.

  \item[L3: Noise Protocol] (Future)
    The Noise protocol framework will provide forward secrecy and
    post-quantum resistance in a future release.
\end{description}

\subsection{DokiLink Sequence Diagram}
\label{sec:dokilink-sequence}

The following sequence diagram shows the establishment of a DokiLink
tunnel between two peers:

\begin{figure}[htbp]
  \centering
  \begin{tikzpicture}[
    node distance=0.4cm,
    every node/.style={font=\sffamily\footnotesize},
    participant/.style={
      rectangle, draw=nodeblue!70, fill=fillblue,
      minimum width=2.5cm, minimum height=0.7cm,
      text=textdark, align=center
    },
    msg/.style={
      -{Stealth[length=2.5mm]}, draw=textdark, line width=0.5pt
    },
    msgRet/.style={
      msg, dashed
    }
  ]
    % Participants
    \node[participant] (peerA) {Peer A};
    \node[participant, right=5cm of peerA] (peerB) {Peer B};

    % Lifelines
    \draw[dashed, linegray] (peerA) -- ++(0,-7);
    \draw[dashed, linegray] (peerB) -- ++(0,-7);

    % Messages
    \draw[msg] ([yshift=-1cm]peerA.south) -- 
      node[above, font=\sffamily\tiny] {1. hello (ed25519 signature)} 
      ([yshift=-1cm]peerB.north);

    \draw[msg] ([yshift=-2cm]peerB.south) -- 
      node[above, font=\sffamily\tiny] {2. peer list (signed)} 
      ([yshift=-2cm]peerA.north);

    \draw[msg] ([yshift=-3cm]peerA.south) -- 
      node[above, font=\sffamily\tiny] {3. TLS ClientHello} 
      ([yshift=-3cm]peerB.north);

    \draw[msgRet] ([yshift=-4cm]peerB.south) -- 
      node[above, font=\sffamily\tiny] {4. TLS ServerHello + cert} 
      ([yshift=-4cm]peerA.north);

    \draw[msg] ([yshift=-5cm]peerA.south) -- 
      node[above, font=\sffamily\tiny] {5. Forward proxy request} 
      ([yshift=-5cm]peerB.north);

    \draw[msgRet] ([yshift=-6cm]peerB.south) -- 
      node[above, font=\sffamily\tiny] {6. Proxied response} 
      ([yshift=-6cm]peerA.north);

    % Self-message for key generation
    \draw[msg] ([yshift=-0.5cm]peerA.south) -- ++(1,0)
      node[right, font=\sffamily\tiny, xshift=2pt] {gen CA + keys}
      -- ++(0,-0.3) -- ++(-1,0);

    % Loop fragment
    \node[draw=linegray!60, dashed, rounded corners=2pt,
      fit=(peerA|-{0,-4.8cm})(peerB|-{0,-5.5cm}),
      inner sep=3pt, fill=white] (loop) {};
    \node[font=\sffamily\tiny, text=textgray, anchor=north west]
      at (loop.north west) {loop: gossip tick};
  \end{tikzpicture}
  \caption{DokiLink-Lite connection establishment sequence.}
  \label{fig:dokilink-sequence}
\end{figure}

\subsection{Peer Discovery}
\label{sec:peer-discovery}

DokiLink supports two discovery mechanisms:

\begin{itemize}[leftmargin=*, itemsep=3pt]
  \item \textbf{Static peers:} Managed via \texttt{doki link add/rm}
        and stored in \texttt{peers.json}. This is the primary
        mechanism for cross-network peers.

  \item \textbf{mDNS discovery:} Available when built with
        \texttt{-tags netlink\_mdns}. Uses multicast DNS to discover
        peers on the local network. Default builds include a no-op
        stub.
\end{itemize}

\subsection{Gossip Protocol}
\label{sec:gossip}

The gossip protocol exchanges ed25519-signed JSON messages over TCP,
capped at 4\,KiB per message. A 15-second tick discovers new peers,
while a 30-second tick refreshes from the static configuration. The
protocol supports three message types: \texttt{hello}, \texttt{peer},
and \texttt{container}.

\section{Rootless Networking (pasta)}
\label{sec:pasta}

For users without root, the pasta utility provides TCP/UDP connectivity
without TAP devices. Doki's \texttt{pkg/network/rootless.go} shells
out to \texttt{pasta} for port forwarding when the standard bridge
network is unavailable.

\section{IPv6 Support}
\label{sec:ipv6}

Doki supports dual-stack IPv4/IPv6 networking. When IPv6 is available
on the host, containers receive both IPv4 and IPv6 addresses from
the bridge network. The DNS server responds to both A (IPv4) and
AAAA (IPv6) queries.

\subsection{Reverse DNS}
\label{sec:reverse-dns}

The \texttt{arpaToIP} function handles reverse DNS lookups for both
address families:

\begin{itemize}[leftmargin=*, itemsep=3pt]
  \item \textbf{IPv4:} \texttt{.in-addr.arpa.} suffix with 4 octets
  \item \textbf{IPv6:} \texttt{.ip6.arpa.} suffix with 32 nibbles
\end{itemize}

This enables proper PTR record resolution for container-to-container
communication in dual-stack environments.

\subsection{Dual-Stack Configuration}
\label{sec:dual-stack}

When both IPv4 and IPv6 are available, the DNS server returns both
A and AAAA records for container hostnames. The client resolver
selects the appropriate address family based on host routing
capabilities.
